\documentclass[10pt,a4paper]{amsart}
\usepackage{amsmath,amssymb,amsthm,graphicx}
% For commutative diagrams we use Xy-pic
\usepackage[all]{xy}

% We use 2cell for 2-commutative diagrams.
\xyoption{2cell}
% Theorem environments.
%
\theoremstyle{plain}
\newtheorem{theorem}[subsection]{Theorem}
\newtheorem{proposition}[subsection]{Proposition}
\newtheorem{lemma}[subsection]{Lemma}

\theoremstyle{definition}
\newtheorem{definition}[subsection]{Definition}
\newtheorem{example}[subsection]{Example}
\newtheorem{exercise}[subsection]{Exercise}
\newtheorem{situation}[subsection]{Situation}

\theoremstyle{remark}
\newtheorem{remark}[subsection]{Remark}
\newtheorem{remarks}[subsection]{Remarks}

\numberwithin{equation}{subsection}

% Macros
%
\def\lim{\mathop{\mathrm{lim}}\nolimits}
\def\colim{\mathop{\mathrm{colim}}\nolimits}
\def\Spec{\mathop{\mathrm{Spec}}}
\def\Hom{\mathop{\mathrm{Hom}}\nolimits}
\def\Ext{\mathop{\mathrm{Ext}}\nolimits}
\def\SheafHom{\mathop{\mathcal{H}\!\mathit{om}}\nolimits}
\def\SheafExt{\mathop{\mathcal{E}\!\mathit{xt}}\nolimits}
\def\Sch{\mathit{Sch}}
\def\Mor{\mathop{\mathrm{Mor}}\nolimits}
\def\Ob{\mathop{\mathrm{Ob}}\nolimits}
\def\Sh{\mathop{\mathit{Sh}}\nolimits}
\def\NL{\mathop{N\!L}\nolimits}
\def\CH{\mathop{\mathrm{CH}}\nolimits}
\def\proetale{{pro\text{-}\acute{e}tale}}
\def\etale{{\acute{e}tale}}
\def\QCoh{\mathit{QCoh}}
\def\Ker{\mathop{\mathrm{Ker}}}
\def\Im{\mathop{\mathrm{Im}}}
\def\Coker{\mathop{\mathrm{Coker}}}
\def\Coim{\mathop{\mathrm{Coim}}}

% Boxtimes
%
\DeclareMathSymbol{\boxtimes}{\mathbin}{AMSa}{"02}

%
% Macros for moduli stacks/spaces
%
\def\QCohstack{\mathcal{QC}\!\mathit{oh}}
\def\Cohstack{\mathcal{C}\!\mathit{oh}}
\def\Spacesstack{\mathcal{S}\!\mathit{paces}}
\def\Quotfunctor{\mathrm{Quot}}
\def\Hilbfunctor{\mathrm{Hilb}}
\def\Curvesstack{\mathcal{C}\!\mathit{urves}}
\def\Polarizedstack{\mathcal{P}\!\mathit{olarized}}
\def\Complexesstack{\mathcal{C}\!\mathit{omplexes}}
% \Pic is the operator that assigns to X its picard group, usage \Pic(X)
% \Picardstack_{X/B} denotes the Picard stack of X over B
% \Picardfunctor_{X/B} denotes the Picard functor of X over B
\def\Pic{\mathop{\mathrm{Pic}}\nolimits}
\def\Picardstack{\mathcal{P}\!\mathit{ic}}
\def\Picardfunctor{\mathrm{Pic}}
\def\Deformationcategory{\mathcal{D}\!\mathit{ef}}
\usepackage{cleveref}
\crefname{originallemma}{original lemma}{original lemmas}
\Crefname{originallemma}{Original lemma}{Original lemmas}
\makeatletter
\newlabel{lemma-tensor-product-locally-constant}{{095D}{source}{}{sourcepointer-1}{}}
\newlabel{lemma-tensor-product-locally-constant@cref}{{[originallemma][0][]095D}{[1][0][]source}{}{}{}}
\makeatother
\title{The tensor product of constructible sheaves of modules on the small etale site is constructible}
\author{The Stacks Project Authors}
\date{stacks-project:etale-cohomology.tex:local-snapshot-20261011; GNU FDL 1.2 or later}
\begin{document}
\maketitle
\noindent\textit{Editorial note.} All delivered mathematical text is the author's own, copied from the pinned Stacks Project chapter file etale-cohomology.tex; no statement or proof step was rewritten, reordered or removed.\par\medskip
\noindent\textit{Editorial note.} The proof is the author's complete argument: the affine reduction, the common refinement and the locally constant case are all inside the counted span.\par\medskip
\noindent\textit{Editorial note.} The label lemma-tensor-product-constructible occurs in no other staged Stacks chapter, so no cross-chapter restatement had to be weighed for this card.\par\medskip
\noindent\textit{Editorial note.} The single labelled reference occurring inside the counted spans is carried in assembly.external\_references as a non-load-bearing navigation pointer to the same chapter.\par\medskip
\noindent\textit{Editorial note.} The preamble is the chapter's own native preamble excerpt already used for this source: the Xy-pic setup (native preamble lines 14-18) and the native theorem declarations with the mathematical macros (native preamble lines 154-217).\par\medskip
\noindent\textit{Allowed standard prerequisites (editorial).}\begin{itemize}
\item Constructible sheaf of Lambda-modules on X\_etale. Same chapter, the definition of constructibility by a partition into locally closed constructible strata on which the sheaf is locally constant. Supplies the objects whose tensor product is shown to be constructible.
\item Common refinement of two partitions into constructible locally closed strata. Same chapter, the standard fact that two such partitions have a common refinement of the same type, used to reduce constructibility statements to a single partition. Supplies the reduction step of the proof.
\item Tensor product of locally constant sheaves is locally constant, and pullbacks commute with tensor product. Same chapter, the corresponding statement for locally constant sheaves and the commutation of pullback with tensor product. Supplies the termwise statement on the common refinement.
\end{itemize}\medskip
\ifdefined\hypertarget\hypertarget{sourcepointer-1}{}\fi
\noindent\textit{Source reference.} Original lemma 095D: sources/stacks-etale-cohomology/pinned-original-source.tex, source lines 10071--10071. Same-chapter navigation pointer invoked at proof line 12001-12002, where the author reduces the statement to the corresponding statement for locally constant sheaves. The pointer only resolves that reference for the reader; the cited lemma is not reproduced inside the counted span, and no bibliography citation is involved.\par\medskip
\begin{lemma}
\label{lemma-tensor-product-constructible}
Let $X$ be a scheme. Let $\Lambda$ be a Noetherian ring.
The tensor product of two constructible sheaves of $\Lambda$-modules
on $X_\etale$ is a constructible sheaf of $\Lambda$-modules.
\end{lemma}

\begin{proof}
The question immediately reduces to the case where $X$ is affine.
Since any two partitions of $X$ with constructible locally
closed strata have a common refinement of the same type and
since pullbacks commute with tensor product we reduce to
Lemma \ref{lemma-tensor-product-locally-constant}.
\end{proof}

\end{document}
